%%
%% REFLEX: ICAIF 2026 submission (ACM sigconf).
%%
%% ############################################################
%% ##  VERSION 2 (CURRENT).  Supersedes archive/main_v1.tex.  ##
%% ############################################################
%%
%% v2 (2026-08-07): rewritten for ICAIF house style. Same results, same
%% figures, same tables, same claim scope as v1; shorter sentences, an
%% explicit contributions list, a separate Related Work section, and no
%% dash-parenthetical asides. Prose budget held below v1 so the paper
%% stays within the 8-page limit. The v1 source and its compiled PDF are
%% kept verbatim under archive/.
%%
%% Build: compiles on Overleaf with the standard acmart bundle
%% (TeX Live >= 2022). Main file: main.tex; bibliography: references.bib;
%% figures under figures/.
%%
%% This is the DE-ANONYMIZED working version (real authors + public repo
%% link). ICAIF review is double-blind: before CMT submission, switch the
%% class line to
%%   \documentclass[sigconf,anonymous,review]{acmart}
%% and swap the Reproducibility footnote for an anonymized mirror.
%% Camera-ready TODOs are marked "TODO(camera-ready)".
%%
\documentclass[sigconf]{acmart}

%% Suppress the ACM reference-format block and copyright box for the
%% pre-acceptance draft. TODO(camera-ready): remove these two lines and
%% insert the rights block ACM supplies on acceptance.
\settopmatter{printacmref=false}
\setcopyright{none}

\acmConference[ICAIF '26]{7th ACM International Conference on AI in
  Finance}{November 14--17, 2026}{Milan, Italy}
\acmYear{2026}
\copyrightyear{2026}

%% acmart already loads amsmath, booktabs, graphicx, hyperref, and the
%% newtxmath fonts (do NOT load amssymb on top of newtxmath).

%% ---------------------------------------------------------------- macros --
\newcommand{\eps}{\varepsilon}
\newcommand{\BR}{\mathrm{BR}}
\newcommand{\Neff}{N_{\mathrm{eff}}}
\newcommand{\hSP}{h_{\mathrm{SP}}}
\newcommand{\hPO}{h_{\mathrm{PO}}}
\newcommand{\gPO}{\gamma_{\mathrm{PO}}}
\newcommand{\geff}{\gamma_{\mathrm{eff}}}
\newcommand{\meff}{m_{\mathrm{eff}}}
\newcommand{\Kdb}{K_{\mathrm{db}}}
\newcommand{\Kmax}{K_{\max}}
\newcommand{\hRef}{h_{\mathrm{ref}}}  % NB: \href is taken by hyperref
\newcommand{\gbar}{\bar{g}}
\newcommand{\Wone}{W_1}
\newcommand{\code}[1]{\texttt{#1}}

\begin{document}

\title[REFLEX: Analytic Stability Boundaries for Performative Market Making]%
{REFLEX: Reflexive Equilibrium Fixed-point Learning for Endogenous
eXchanges}
\subtitle{Analytic Stability Boundaries for Performative Market Making in
OTC Corporate Bond Markets}

%% -------------------------------------------------------------- authors --
%% Co-first authors (equal contribution); Vignesh is corresponding (the
%% first \authornote mark renders as the asterisk).
%% TODO(camera-ready): add \orcid{} for both authors (ICAIF requires ORCID
%% iDs at camera-ready).
\author{Vignesh Nagarajan}
\authornote{Corresponding author.}
\authornote{Both authors contributed equally to this research.}
\affiliation{%
  \institution{Texas A\&M University}
  \department{Brown Foundation Scholar}
  \city{College Station}
  \state{TX}
  \country{USA}}
\email{vigneshn26@tamu.edu}

\author{Shriraghav Ashok}
\authornotemark[2]
\affiliation{%
  \institution{University of California, Berkeley}
  \department{Haas School of Business}
  \city{Berkeley}
  \state{CA}
  \country{USA}}
\email{sashok24@berkeley.edu}

%% ------------------------------------------------------------- abstract --
\begin{abstract}
A machine-learned bond dealer reshapes the data it is trained on. Tighter
quotes win volume and also attract informed flow that trades ahead of
adverse price moves, so retraining on the induced flow closes a feedback
loop. Performative prediction bounds when that loop converges, but its
condition is stated in Lipschitz constants of an unspecified loss that no
desk can evaluate before deploying. We present \textbf{REFLEX}, which pins
those constants to the microstructure primitives of a structural
over-the-counter corporate-bond market-making model: fill-curve curvature,
P\&L scale, and the slope of informed flow against spread. Stability becomes
computable before deployment, and the boundary extends in closed form to
gradient correction, competing dealers, finite-sample certificates,
128-bond factor universes, and lazy retraining cadences. In a learned
simulator loop, measured contraction moduli track the closed form within
$8\%$, and competition amplifies instability by the predicted effective
dealer count. A structurally anchored corrected loop converges where blind
retraining collapses and where a free-form learned correction fails.
Calibrated on 36 years of public market data, the boundary yields a daily
fragility index whose stability headroom collapses roughly $4.4\times$ from
calm to crisis, and all 66 load-bearing identities are machine-checked.
REFLEX gives desks and supervisors a stability margin for algorithmic
liquidity provision that is computable before a quoting model goes live.
\end{abstract}

%% ------------------------------------------------------ CCS and keywords --
\begin{CCSXML}
<ccs2012>
<concept>
<concept_id>10010147.10010257</concept_id>
<concept_desc>Computing methodologies~Machine learning</concept_desc>
<concept_significance>500</concept_significance>
</concept>
<concept>
<concept_id>10010405.10010481.10010485</concept_id>
<concept_desc>Applied computing~Economics</concept_desc>
<concept_significance>500</concept_significance>
</concept>
<concept>
<concept_id>10003752.10003790.10011742</concept_id>
<concept_desc>Theory of computation~Market equilibria</concept_desc>
<concept_significance>300</concept_significance>
</concept>
</ccs2012>
\end{CCSXML}
\ccsdesc[500]{Computing methodologies~Machine learning}
\ccsdesc[500]{Applied computing~Economics}
\ccsdesc[300]{Theory of computation~Market equilibria}

\keywords{performative prediction, market making, OTC corporate bonds,
retraining stability, systemic risk, machine-checked proofs}

\maketitle

%% ========================================================================
\section{Introduction}

Corporate bonds trade over the counter. A client requests a quote, a dealer
posts a half-spread $h$, and the dealer's inventory absorbs the flow.
Quoting on these desks is increasingly algorithmic, and quoting policies are
increasingly fit to historical flow.

That flow is not exogenous. A tighter quote wins more volume, and it also
summons more informed (``toxic'') flow that picks the dealer off before
adverse price moves. A wider quote starves both channels. The policy $\phi$
therefore induces the distribution $D(\phi)$ it will next be trained on.
Fitting on the induced flow, redeploying, and refitting is exactly
\emph{repeated risk minimization} (RRM) in the sense of performative
prediction \cite{perdomo2020performative}. RRM converges to a
performatively stable point if and only if $\eps<\gamma/\beta$, where $\eps$
bounds the sensitivity of the distribution map and $\gamma,\beta$ are the
strong convexity and smoothness of the loss.

The theorem is sharp, but its constants are Lipschitz \emph{assumptions}. A
desk cannot evaluate $\eps$, $\beta$, or $\gamma$ before deploying, so it
cannot know whether its own retraining cadence is stable. A supervisor
watching many dealers retrain against one shared pool of informed flow
cannot know whether competition among learners is manufacturing systemic
fragility. Post-trade TRACE reporting reveals the spreads that result, not
the retraining loop that produced them. Both questions are quantitative and
both need the constants rather than the abstraction. The object we compute
is a market-level stability margin for algorithmic liquidity provision,
which places the problem in the AI-governance and systemic-risk scope of AI
in finance.

REFLEX (\emph{Reflexive Equilibrium Fixed-point Learning for Endogenous
eXchanges}) realizes performative prediction \emph{inside} a structural OTC
market-making model
in the Gu\'eant--Lehalle--Fern\'andez-Tapia (GLFT) family. It derives the
loop constants from microstructure primitives, then verifies every closed
form against a learned simulator loop. We contribute six results, each
derived and then confirmed or falsified by experiment:
\begin{enumerate}
\item[\textbf{R1}] \emph{An analytic stability boundary.} We give $\gamma$,
  $\beta$, and $\eps$ in closed form from fill-curve curvature, P\&L scale,
  and toxic-flow slope, so the modulus $m=\eps\beta/\gamma$ is known before
  deployment. Measured moduli track it within $8\%$ where the loop
  contracts.
\item[\textbf{R2}] \emph{A performative-gradient (PerfGD) correction},
  closed-form and estimated. One extra scalar per step converges past the
  RRM boundary, and its structurally anchored learned counterpart settles at
  the realized performative optimum where blind retraining collapses.
\item[\textbf{R3}] \emph{A multi-dealer boundary} $\eps<\gamma/(\Neff\beta)$.
  Competition destabilizes the market a factor $\Neff$ before any single
  dealer would, and a genuine shared-pool market measures $1.74\times$ and
  $3.16\times$ amplification at $N=2,3$.
\item[\textbf{R4}] \emph{Finite-sample robust certificates.} An
  $O(1/\sqrt{n})$ ambiguity radius bought by common random numbers (CRN)
  issues stable, unstable, or \emph{undecided} verdicts, separating
  statistical from structural uncertainty.
\item[\textbf{R5}] \emph{Factor-model scaling.} The $d{\times}d$ modulus
  matrix has spectral boundary $\rho(M)<1$ at $O(dk^2)$ cost, and
  $\rho(M)$ is flat from $8$ to $128$ bonds on our calibration.
\item[\textbf{R6}] \emph{Lazy deployment.} The $K$-step retraining map
  $\mu(K)=-m+c^K(1+m)$ yields deadbeat and maximal-stability cadences in
  closed form. Both parameter-free predictions land on the measured probe.
\end{enumerate}

The system side contributes an \emph{un-blinded} learned market-response
operator whose training identifies $dD/d\phi$; three independent instruments
for $\eps$, namely a CRN best-response probe, an entropic optimal-transport
estimate \cite{cuturi2013sinkhorn,feydy2019interpolating,
peyre2019computational}, and a fitted informed-flow curve
\cite{cont2014price}; calibration of the simulator to 36 years of public
market data \cite{dickerson2023priced,friewald2012illiquidity}; and a
verification layer of 66 numerical proof certificates plus Lean~4
\cite{moura2021lean,mathlib2020} formal skeletons. We report negative
results and proxy-level data provenance as such throughout.

%% ========================================================================
\section{Related Work}
\label{sec:related}

Two mature literatures approach the loop from opposite ends. The
performative-prediction literature has sharpened the loop itself:
convergence of RRM and its stochastic variants
\cite{perdomo2020performative,mendlerdunner2020stochastic}, gradient methods
that estimate the distribution response $dD/d\phi$ and optimize performative
risk directly \cite{izzo2021learn,miller2021echo}, decision-dependent
stochastic optimization \cite{drusvyatskiy2023stochastic}, state-dependent
extensions in which the distribution carries its own dynamics
\cite{li2022state,brown2022stateful}, bandit feedback
\cite{jagadeesan2022regret}, multi-agent decision-dependent games
\cite{narang2022multiplayer}, and distributionally robust formulations; see
\cite{hardt2023performative} for a survey. Throughout,
$(\eps,\beta,\gamma)$ remain abstract constants of an unspecified loss.

The market-making literature supplies the missing structure: closed-form
optimal quotes under exponential fill intensities
\cite{avellaneda2008highfrequency,gueant2013dealing}, multi-asset
dimensionality reduction \cite{bergault2021size}, adverse selection and
price reading \cite{barzykin2025adverse}, the stochastic-control toolkit
\cite{cartea2015algorithmic}, and the empirical anatomy of flow--price
interaction \cite{cont2014price}. These models treat the flow response to
quoting as a fixed input. The learning loop that re-estimates the market
from its own induced data is outside their scope.

To our knowledge, no prior work computes the performative constants from
market structure, which is what makes the stability theorem operational.

%% ========================================================================
\section{Market Model and the Retraining Loop}
\label{sec:model}

One deployment is one period of quoting by a dealer posting half-spread $h$.
Skew is zero for the scalar theory; the multi-bond lift is R5. With $\rho$
the latent liquidity ratio, benign notional follows the GLFT exponential
fill curve and informed notional a spread-gated channel:
\begin{equation}
U(h)=A\rho\,e^{-kh},\qquad
\tau(h)=\rho\,\gbar\big(I_b+\alpha f I\,e^{-c_t h}\big),
\label{eq:flows}
\end{equation}
where $A,k$ are the arrival rate and demand elasticity, $I_b,I$ the base and
feedback informed intensities, $\alpha$ adversarial informativeness, $f$ the
toxicity-feedback gain (our control variable), $c_t$ the toxic spread decay,
and $\gbar\in(0,1)$ the mean of the informed traders' $\tanh$ signal gate.
Informed flow trades with the signal, so each unit costs the dealer an
adverse-selection severity $\psi>0$, a state constant from the same Gaussian
gate integrals. The expected one-deployment objective, with P\&L scale $P$,
quoting-cost weight $w$, anchor $\hRef$, and inventory-variance curvature
$\lambda_q\ge0$, is
\begin{equation}
J(h;T)=P\big[\,hA\rho e^{-kh}+hT-\psi T-w(h-\hRef)^2-\lambda_q\,\big].
\label{eq:objective}
\end{equation}
The structural content of ``the learner is blind'' is the frozen toxic level
$T$. Within a deployment the fitted model conditions on the deployed regime,
so the dealer best-responds to $T=\tau(h_{\mathrm{dep}})$ while only the
benign channel varies with its candidate $h$. Retraining is then the cobweb
\begin{equation}
h_{t+1}=\BR(h_t):=\arg\max_h J\big(h;\tau(h_t)\big),\qquad
m:=\big|\BR'(h^*)\big|,
\label{eq:cobweb}
\end{equation}
which contracts to the performatively stable spread $\hSP$ iff $m<1$. The
simulator realizes \eqref{eq:flows} to \eqref{eq:cobweb} with multi-bond
flow, an informed-flow saturation cap, and a liquidity field that
\emph{inflates with realized flow}. Those endogenous channels are omitted by
the frozen closed forms and made visible by the measurements below.

%% ========================================================================
\section{Closed-Form Stability Theory}
\label{sec:theory}

Full derivations, with proofs, regularity conditions, and saturated
corollaries, are documents D1--D6 in the repository (see the
\emph{Reproducibility} statement). We state the objects the experiments
test.

\textbf{R1 (analytic boundary).} Differentiating \eqref{eq:objective}, the
three Perdomo constants are closed forms in the configuration:
\begin{align}
\gamma&=P\big[2w+A\rho k e^{-kh}(2-kh)+\lambda_q\big],\qquad \beta=P,
\notag\\
\eps(h)&=\rho\,\gbar\,\alpha f I c_t\,e^{-c_t h}.
\label{eq:constants}
\end{align}
Curvature is fill-curve curvature plus quoting-cost stiffness, smoothness is
the P\&L scale, and $\eps$ is the toxic-flow slope $|d\tau/dh|$. The
implicit-function theorem gives $\BR'=-\eps\beta/\gamma$, an oscillatory
cobweb, hence
\begin{equation}
m(h^*)=\frac{\eps\beta}{\gamma}
=\frac{\rho\gbar\alpha f I c_t\,e^{-c_t h^*}}
       {2w+A\rho k\,e^{-k h^*}(2-k h^*)+\lambda_q},
\label{eq:modulus}
\end{equation}
with $P$ cancelling and $h^*$ the self-consistent fixed point of
\eqref{eq:cobweb}. The loop is stable iff $\eps<\gamma/\beta$, i.e.\ $m<1$.
Every term is evaluable before any loop is run, so the boundary is a
\emph{prediction}. Two structural corollaries matter for honest measurement.
First, adverse selection ($\psi T$) is constant in $h$ within a deployment,
so it sets the profit level (the echo-chamber gap of R2) but drops out of
the boundary. Second, at the self-consistent fixed point the modulus
\emph{saturates}: widening spreads decay $\eps(h^*)$ exponentially, so at
default-like constants the fixed-point $m$ never crosses $1$. This is
defensive widening. Measured instability is therefore a statement about the
retraining map at the \emph{operating} spread, and we report it as such.

\textbf{R2 (PerfGD correction).} The performative objective is
$\Phi(h)=J(h;\tau(h))$. Un-blinding adds one closed-form term to the blind
gradient $G(h)=\partial_1 J(h;\tau(h))$:
\begin{equation}
\Phi'(h)=G(h)+\Delta(h),\qquad
\Delta(h)=-\beta\,(h-\psi)\,\eps(h),
\label{eq:perfgd}
\end{equation}
where $\Delta=\partial_T J\cdot(d\tau/dh)$ is the distribution-response
term, supplied entirely by the R1 closed forms. The corrected ascent costs
one scalar more per step than RRM \cite{izzo2021learn}. Its convergence is
governed not by the cobweb modulus but by the objective curvature at the
performative optimum $\hPO$,
\begin{equation}
\gPO=\gamma+\beta\eps\,(2+c_t\psi-c_t\hPO),
\label{eq:gpo}
\end{equation}
which stays positive where $m>1$, so the corrected loop converges beyond the
RRM boundary. The stable-versus-optimal separation, the ``echo-chamber
gap'' of \cite{miller2021echo}, is
$\hSP-\hPO=\beta\eps(\hSP-\psi)/(\gamma+\beta\eps)=O(\eps)$ in decision
space and $O(\eps^2)$ in value.

\textbf{R3 (multi-dealer systemic risk).} $N$ dealers share one informed
pool with toxic spillover $\kappa\in[0,1]$, so flow displaced by one
dealer's widening lands on the others. The joint best-response Jacobian is a
rank-one common-mode coupling. Its differential eigenvalues carry
$(1-\kappa)m_1$ and its common mode $-m_1\Neff$ with
$\Neff=1+\kappa(N-1)$, so the joint loop is stable iff
\begin{equation}
\eps<\frac{\gamma}{\Neff\,\beta},
\qquad
N_c=1/m_1
\label{eq:multidealer}
\end{equation}
is the critical dealer count at which a market of individually stable
dealers ($m_1<1$) turns systemically unstable. Competition manufactures a
synchronized cobweb, making fragility a market property rather than a desk
property.

\textbf{R4 (robust boundary).} With $n$ probe estimates of $\eps$ (sample
mean $\hat\eps_n$, std $s$), the certificate
\begin{equation}
\hat\eps_n+\delta_n<\frac{\gamma}{\beta},\qquad
\delta_n=\frac{z_{1-a}\,s}{\sqrt{n}},
\label{eq:robust}
\end{equation}
declares \emph{stable}; $\hat\eps_n-\delta_n>\gamma/\beta$ declares
\emph{unstable}; otherwise \emph{undecided}. The parametric $O(1/\sqrt{n})$
radius is bought by the CRN pairing of the probe, since a naive finite
difference of independent runs achieves only $O(n^{-1/3})$ after
bias--variance balancing. Pinning a crossing needs
$n_{\mathrm{req}}=O(\Delta^{-2})$ in the distance $\Delta$ to the boundary,
so statistical uncertainty is separated from structural uncertainty by
construction. A distribution-free quantile calibration of $\delta_n$, new
here, shows $z\,s$ is conservative except under contamination
(Sec.~\ref{sec:tuning}).

\textbf{R5 (factor scaling).} For $d$ bonds, stability is governed by the
modulus matrix $M=\beta\,\Gamma^{-1}E$ with $E=\mathrm{diag}(\eps_i)$ and
$\Gamma$ the cross-bond curvature, and the loop contracts iff $\rho(M)<1$. A
$k$-factor covariance makes $\Gamma$ diagonal-plus-low-rank, so $\rho(M)$
costs $O(dk^2)$ by Woodbury, with truncation error linear in the residual
factor variance $\lambda_{k+1}(C)$.

\textbf{R6 (lazy deployment).} Suppose each deployment takes only $K$ inner
gradient steps on the frozen-$T$ objective (step $\eta_h$, contraction
$c=1-\eta_h\gamma$), warm-started from the deployed spread, instead of the
exact best response. The outer map slope becomes
\begin{equation}
\mu(K)=-m+c^K(1+m),
\label{eq:lazy}
\end{equation}
a one-parameter interpolation from inertia ($\mu(0)=1$) to the exact cobweb
($\mu\to-m$). Three consequences follow in closed form. A \emph{deadbeat}
cadence $\Kdb=\ln\!\big(m/(1+m)\big)/\ln c$ gives $\mu=0$ and one-shot
convergence. For $m>1$ a stability window $K\le\Kmax=
\ln\!\big((m-1)/(m+1)\big)/\ln c$ keeps an RRM-unstable market stable, which
is the quantitative form of the greedy-versus-lazy gap in
\cite{perdomo2020performative,mendlerdunner2020stochastic}. And an observer
who fits the plain boundary to a lazy loop reads a two-branch
\emph{effective} curvature $\geff(K)=\gamma\,m/|\mu(K)|$, under-estimating
stiffness below the equal-modulus cadence and over-estimating it above.

%% ========================================================================
\section{Experimental Setup}
\label{sec:system}

\subsection{Learned Operator and the Four Loop Modes}
The simulator is a multi-bond OTC market with benign and informed clients,
inventory, an informed-flow saturation cap, and a latent liquidity field
inflated by realized flow. The learned market-response operator $T_\theta$
(MLP heads over per-bond features) predicts next-deployment flow
distributions conditioned on a per-deployment \emph{policy summary}, and is
fit over a sliding window of past deployments. With window $\ge2$ the
derivative of its prediction w.r.t.\ the summary, the learned $dD/d\phi$, is
identified by automatic differentiation; the single-window fit is
structurally blind and gives the RRM baseline. Four retraining modes close
the loop: \textbf{(i)}~blind RRM; \textbf{(ii)}~\emph{PerfGD-analytic},
using the closed-form $\Delta$ of \eqref{eq:perfgd} as a surrogate gradient;
\textbf{(iii)}~\emph{PerfGD-learned}, consuming the free-form learned
$dD/d\phi$, reported as a negative result; and
\textbf{(iv)}~\emph{PerfGD-structural}, which never asks the network for a
derivative. The structural mode fits the theory's own response families
\begin{equation}
\hat\tau(h)=\hat C_0+\hat C_1 e^{-\hat c h},\qquad
\hat u(h)=\hat A_u e^{-\hat k_u h},
\label{eq:structfit}
\end{equation}
plus the realized severity $\hat\psi$, to the loop's \emph{own} deployment
history (per-bond, per-step samples over a 12-deployment window; a small
quote jitter provides within-deployment identification that the echo chamber
cannot destroy). It then ascends $\hat\Phi'=\hat G+\hat\Delta$ with step
$1/\hat\gPO$ under a trust region ($35\%$ relative step cap) and an
anti-echo freeze that refuses to steer on unidentified fits. Each safeguard
encodes a measured failure. Exponential response fits on a narrow spread
window are line-degenerate and extrapolate badly, and the unconstrained
early ascent rails against its spread cap.

\subsection{Three Instruments for $\eps$ and the Probe Protocol}
\label{sec:instruments}
(a)~\emph{CRN best-response probe}: paired deployments at $\hRef\pm\delta$
under common random numbers. The finite difference of best responses
estimates $m$ (and $\eps=\hat m\gamma/\beta$), and a signed $K$-step variant
measures $\mu(K)$ for R6. (b)~\emph{Optimal transport}:
$\hat\eps=\Wone\big(D(h{+}\delta),D(h{-}\delta)\big)/2\delta$ is the exact
Wasserstein sensitivity of performative prediction, computed via debiased
log-domain Sinkhorn divergences
\cite{cuturi2013sinkhorn,feydy2019interpolating} with the blur set
\emph{scale-relatively} ($0.02\times$ pooled std, tuned against the exact
1-D quantile $\Wone$; Sec.~\ref{sec:tuning}). (c)~\emph{Fitted
informed-flow curve}: bin realized toxic notional against deployed spread,
fit \eqref{eq:flows}, differentiate \cite{cont2014price}. A pre-run audit
showed three protocol facts that matter at first order: probe at the
operating spread rather than the analytic fixed point; keep collection
jitter at $0.05$, since at $0.2$ the BR probe inflates ${\sim}3\times$; and
report medians and interquartile ranges (IQR) over 8 seeds with the R4
robust bands at every grid point. We sweep the gain $f$ rather than
adversariality $\alpha$. High $\alpha$ drives the dealer to wide spreads
where $e^{-c_t h}$ has decayed, so the $\alpha$-response is non-monotone
(measured hump $0.08\to1.83\to0.67$). That confound is the quantitative case
for $f$ as the control variable.

\subsection{Real-Data Calibration and Provenance}
\label{sec:data}
The simulator is calibrated per (rating $\times$ volatility regime) on a
public, verified panel: ${\sim}36$ years of daily and ${\sim}70$ years of
monthly series (CBOE VIX as the $\sigma$ proxy and regime classifier; WTI;
the Fed H.15 10-year; Shiller equity data; gold/CPI), TRACE-derived
corporate-bond factors of \cite{dickerson2023priced} whose liquidity risk
factor is the primary $\eps$ proxy, and monthly returns for 212 real-CUSIP
bonds. Half-spread proxies combine VIX-implied spreads with the illiquidity
add-on of \cite{friewald2012illiquidity,bao2011illiquidity}, and
$\lambda(h)=Ae^{-kh}$ intensities are fit per cell by maximum likelihood.
\emph{Provenance, stated plainly:} this is not trade-level TRACE.
Dealer-side prints, per-dealer inventories, and per-bond $(A,k)$ would
require TRACE Enhanced. Only $(A,k,\sigma,h)$ are data-identified here, the
toxic channel is structurally scaled by documented ratios, and the
crisis-regime intensity fit is degenerate ($k=0$, $n=74$ days) and flagged
wherever it binds. Regime \emph{ordering} of the boundary is data-driven;
absolute critical gains are not. Calibrated configurations run in
per-\$100-par units, and every probe width, tolerance, and OT blur is
scale-relative by convention.

\subsection{Verification Layer}
\label{sec:verification}
Every load-bearing identity, inequality, and dynamical claim of R1--R6 is
re-derived numerically by 66 \emph{proof certificates}, run on both the raw
and the calibrated real-unit configurations: central differences of the
\emph{composed} best-response map against the assembled $-\eps\beta/\gamma$
with no shared code path, eigensolves against spectral formulas,
Monte-Carlo rate fits against claimed exponents, and the actual dynamics run
rather than assumed. Building the layer caught real errors before the
paper-grade run: a convention mismatch in which the 1-D frozen-gradient
helpers omit the inventory curvature, so their slopes are governed by
$\gamma-P\lambda_q$; and two falsified first drafts, an ``always-stiffer''
$\geff$ claim corrected to the two-branch law of R6, and a mis-signed
heavy-tail prediction for the R4 radius. The logical skeletons
of all six results are additionally formalized in Lean~4 against mathlib
\cite{moura2021lean,mathlib2020}. These are reviewed statements whose
compilation is pending a toolchain, so the numerical certificates are the
verification of record. The full suite (152 tests, 11 experiments) runs
CPU-only in ${\sim}25$ minutes, deterministic from (configuration, seed).

%% ========================================================================
\section{Results}
\label{sec:results}

\subsection{Machine-Checked Identities}
Table~\ref{tab:certs} summarizes the certificates: 66/66 pass, with worst
residuals far inside tolerance. The one excluded cell is deliberate. The R2
beyond-boundary \emph{demo dynamics} certificate pins absolute raw-unit
constants, and imposing those on per-\$100-par calibrated units is exactly
the class of unit bug the conventions forbid. It excludes itself, and the
calibrated pass proves the identities in real units.

\begin{table}[t]
\caption{Numerical proof certificates (raw + calibrated real-unit
configurations). Each family's worst residual sits inside its stated
tolerance.}
\label{tab:certs}
\small
\begin{tabular}{@{}lccrr@{}}
\toprule
Certificate family & raw & calib. & worst resid. & tol. \\
\midrule
R1 analytic boundary        & 5/5 & 5/5 & $6.8{\times}10^{-2}$ & $8{\times}10^{-2}$ \\
R2 PerfGD identities        & 4/4 & 4/4 & $1.3{\times}10^{-6}$ & $2{\times}10^{-2}$ \\
R2 PerfGD demo dynamics     & 4/4 & --  & $1.4{\times}10^{-11}$ & $1{\times}10^{-3}$ \\
R3 multi-dealer             & 6/6 & 6/6 & $2.2{\times}10^{-16}$ & $1{\times}10^{-9}$ \\
R4 robust boundary          & 7/7 & 7/7 & $3.2{\times}10^{-2}$ & $1{\times}10^{-1}$ \\
R5 factor scaling           & 3/3 & 3/3 & $0.0$                & $1{\times}10^{-9}$ \\
R6 lazy deployment          & 6/6 & 6/6 & $1.1{\times}10^{-16}$ & $1{\times}10^{-12}$ \\
\bottomrule
\end{tabular}
\end{table}

\subsection{A Real-Data Fragility Index (1990--2026)}
\label{sec:fragility}
Because R1 is a closed form of observables, it can be evaluated on every
trading day of the panel. Figure~\ref{fig:fragility} plots the resulting
index, $\mathrm{fragility}(t)=\mathrm{median}(\eps^*)/\eps^*(t)$, and
Table~\ref{tab:fragility} gives its per-regime summary. Three facts survive
all caveats. The stability headroom $\eps^*=\gamma/\beta$ collapses
${\sim}4.4\times$ (investment grade, IG) and ${\sim}4.3\times$ (high yield,
HY) from calm to crisis. HY sits more than $10\times$ below IG in
\emph{every} regime, which is a level effect rather than a crisis effect.
And the modulus at \emph{observed} spreads falls into crisis (IG
$0.85\to0.14$), so dealers widen faster than the toxic channel steepens.
Defensive widening is visible in real data. The index saturates at a crisis
plateau because of the degenerate crisis fit: the GFC and the March-2020
COVID freeze are flagged at the same ceiling, and intra-crisis ranking is
explicitly not identified. Per-cell a-priori boundaries span
$\eps^*=465\to2.2$ from IG-calm to HY-crisis while the fixed-point modulus
stays regime-invariant (${\approx}0.066$) by construction, so the regime
story lives entirely in the headroom.

\begin{figure}[t]
\centering
\includegraphics[width=\columnwidth]{figures/fragility_index.png}
\caption{The daily market-fragility index, 1990--2026: the R1 closed
forms evaluated on real data. Top:
$\mathrm{fragility}(t)=\mathrm{median}(\eps^*)/\eps^*(t)$ for IG (solid)
and HY (dashed), with the Lehman and COVID-freeze dates marked. Bottom:
the VIX spine, colored by volatility regime. The crisis plateau reflects
the degenerate crisis-cell fit and is flagged, not hidden.}
\label{fig:fragility}
\end{figure}

\begin{table}[t]
\caption{Per-regime stability headroom $\eps^*=\gamma/\beta$ and modulus
at observed spreads, evaluated on the 1990--2026 panel.}
\label{tab:fragility}
\small
\begin{tabular}{@{}lrrrr@{}}
\toprule
regime & $\eps^*_{\mathrm{IG}}$ & $\eps^*_{\mathrm{HY}}$ &
$m_{\mathrm{IG}}(h_{\mathrm{obs}})$ & $m_{\mathrm{HY}}(h_{\mathrm{obs}})$\\
\midrule
calm     & 1207.7 & 93.5 & 0.847 & 0.592 \\
normal   &  739.2 & 61.3 & 0.870 & 0.563 \\
elevated &  594.4 & 45.9 & 0.520 & 0.367 \\
stress   &  477.3 & 35.6 & 0.223 & 0.163 \\
crisis   &  275.9 & 21.8 & 0.139 & 0.091 \\
\bottomrule
\end{tabular}
\end{table}

\subsection{Predict-Then-Verify: the Boundary Crossing}
\label{sec:sweep}
Figure~\ref{fig:sweep} is the headline falsification test. We sweep the gain
$f$ (7 values $\times$ 8 seeds), overlay the a-priori curve
\eqref{eq:modulus}, and measure the modulus with the CRN probe. In the
contracting regime the agreement is quantitative: measured median $0.390$
against predicted $0.426$ at $f=2$, an $8\%$ gap. The measured crossing sits
at $f^*\approx3.17$ against the a-priori $4.70$, and that gap is not free.
The triangulation (Table~\ref{tab:triangulation}) independently measures the
deployment's own flow inflating the liquidity field to $\rho\approx2.3$, and
the closed form evaluated at that \emph{realized} state predicts
$f^*\approx2.8$--$3.0$, bracketing the measurement. The R4 certificates
grade the grid exactly as the seed bands warrant: stable through $f=2$,
unstable at $f=4$, undecided where beyond-boundary readings scatter. At
$f=8$ the eight seeds span $0.03$--$2.05$, because past the boundary the
probe is a finite-difference diagnostic rather than a local slope. The
fixed-point curve itself saturates at $0.61$, so the measured instability is
a property of the retraining map at the operating spread, exactly as R1
predicts.

The two distribution-space instruments agree with each other within $14\%$
and sit $2.3$--$2.7\times$ above the realized-state closed form. The
residual is the state-feedback channel that any frozen-state closed form
omits by construction, so the analytic boundary is a \emph{lower anchor}
rather than an unbiased point prediction. That is the operationally safe
direction for a stability certificate.

\begin{figure}[t]
\centering
\includegraphics[width=\columnwidth]{figures/phase_diagram_toxicity_feedback.png}
\caption{Predict-then-verify phase diagram over the feedback gain $f$:
measured modulus (median, IQR band, and R4 robust band over 8 seeds)
against the a-priori closed form at the probe spread and the saturating
fixed-point curve, with the measured ($f^*\approx3.17$) and predicted
($4.70$) crossings marked.}
\label{fig:sweep}
\end{figure}

\begin{table}[t]
\caption{Three-way $\eps$ triangulation at the operating spread, against
the closed form at the a-priori (A2) and at the realized deployment state
($\rho=2.32$, $|g|=0.446$).}
\label{tab:triangulation}
\small
\begin{tabular}{@{}lrr@{}}
\toprule
instrument & $\hat\eps$ & vs.\ realized \\
\midrule
closed form, a-priori state ($\rho=1$) & 0.764 & $0.44\times$ \\
closed form, \textbf{realized state}   & 1.726 & $1.00$ \\
CRN best-response slope                & 0.508 & $0.29\times$ \\
Sinkhorn/$\Wone$                       & 4.578 & $2.65\times$ \\
Fitted informed-flow curve             & 4.011 & $2.32\times$ \\
\bottomrule
\end{tabular}
\end{table}

\subsection{Competition Amplifies Instability (R3)}
\label{sec:dealers}
We probe a \emph{genuine} $N$-dealer market sharing one informed pool, not
an analytic shortcut; the environment reduces bit-for-bit to the
single-dealer market at $N=1$. CRN joint-mode probes measure the
common-mode slope against the predicted $\Neff m_1$
(Table~\ref{tab:dealers}). Amplification is $1.74\times$ at $N{=}2$ and
$3.16\times$ at $N{=}3$ against predicted $2$ and $3$, within $13\%$ and
$5\%$. The differential mode is dead at full spillover
($(1-\kappa)m_1=0$), so the instability is purely the synchronized systemic
channel. The analytic surface (Figure~\ref{fig:dealersurface}) puts the
critical dealer count at $N_c\approx7.9$ for the default single-dealer
modulus, and the simulated three-dealer cobweb already oscillates against
its cap. The governance reading is direct: each desk can satisfy its own
stability check while the market as a whole violates
\eqref{eq:multidealer}, so fragility must be certified at the market level.

\begin{figure}[t]
\centering
\includegraphics[width=\columnwidth]{figures/dealer_phase_diagram.png}
\caption{The R3 systemic surface: joint modulus $m_N=\Neff\,m_1$ over the
$(f,N)$ grid at full spillover ($\kappa=1$), with the $m_N=1$ boundary
contour. Gains that are comfortably stable for a single dealer cross the
systemic boundary as the dealer count grows; the probes of
Table~\ref{tab:dealers} measure the first rows of this surface on the
simulated market.}
\label{fig:dealersurface}
\end{figure}

\begin{table}[t]
\caption{Measured common-mode retraining slope on the shared-pool market
versus the R3 prediction $\Neff m_1$ (full spillover, $\kappa=1$).}
\label{tab:dealers}
\small
\begin{tabular}{@{}crrr@{}}
\toprule
$N$ & measured & predicted & differential mode \\
\midrule
1 & 0.786 & 0.786 & -- \\
2 & 1.369 & 1.571 & 0.0034 \\
3 & 2.480 & 2.357 & 0.0032 \\
\bottomrule
\end{tabular}
\end{table}

\subsection{Closing the Loop-Level Gap: Structural PerfGD (R2)}
\label{sec:gap}
The demo regime is genuinely RRM-unstable ($f=6$, slow toxic decay, cobweb
modulus $1.21$). There the blind cobweb provably diverges while the
corrected 1-D ascent converges, in closed form. The question for the learned
system is whether a \emph{learned} loop can realize the correction.
Figure~\ref{fig:loops} and Table~\ref{tab:loops} answer it from a common
seed. Blind RRM collapses into the echo chamber, ending at half-spread
$0.16$. PerfGD-analytic and the free-form PerfGD-learned mode both fail to
settle. Each consumes the operator's implied objective gradient, and the
run's seam diagnostics show the free-form learned toxic slope starting
right-signed ($-0.84$ against analytic $-1.79$), then flipping sign and
hovering near zero. Capacity is not the problem; identification off the
deployed regime is. The structural mode never asks the network for a
derivative. Fitting \eqref{eq:structfit} to its own deployment history, it
climbs under the trust region and \emph{converges to the realized
performative optimum} at $h=3.193$, with late steps falling
$0.35\to0.02$. That is within $0.7\%$ of its own running estimate
$\hat h_{\mathrm{PO}}=3.172$, and its fitted slope tracks the analytic shape
($-1.33$ against $-2.00$ at the final iterate) while strengthening as the
fit window widens.

\emph{Benchmark honesty.} The frozen-reference closed-form optimum
($\hPO=1.641$) is \emph{not} the target. In this high-intensity regime the
realized market differs from the frozen closed forms at first order through
channels they omit by construction: informed-flow saturation, the
$\rho\approx2.3$ liquidity inflation of Sec.~\ref{sec:sweep}, and severity
drift. The loop is therefore verified against \emph{independent structural
fits} on fresh controlled deployments spanning the operating range. The
settle point (i) zeroes their corrected gradient (residual $<30\%$ of its
start value), (ii) sits strictly inside their blind stable point
($<0.85\times\hat h_{\mathrm{SP}}$, the realized echo-chamber gap, closed),
and (iii) brackets their independent optimum estimate. The scoped claim is
this: un-blinding is proven in closed form, and the structurally anchored
learned loop realizes it at loop level against the realized-market
benchmark, while the free-form learned correction remains a documented
negative result. Anchoring, not capacity, closes the gap.

\begin{figure*}[t]
\centering
\includegraphics[width=\textwidth]{figures/perfgd_ml_loops.png}
\caption{Left: half-spread trajectories of the four learned retraining
loops from a common seed in the RRM-unstable regime (the frozen-reference
$\hPO$ and $\hSP$ lines are context, not the benchmark;
Sec.~\ref{sec:gap}). Blind RRM collapses; the structurally anchored loop
converges to the realized performative optimum. Right: the three-way
ML$\leftrightarrow$theory seam, showing the free-form learned, structurally
fitted, and analytic toxic slopes logged at every deployment. The
free-form slope flips sign and hovers near zero; the structural fit
tracks the analytic shape.}
\label{fig:loops}
\end{figure*}

\begin{table}[t]
\caption{Loop outcomes in the RRM-unstable demo regime (10 deployments,
common seed; convergence verified against independent structural fits).}
\label{tab:loops}
\small
\begin{tabular}{@{}lrl@{}}
\toprule
mode & final $h$ & outcome \\
\midrule
blind RRM          & 0.16 & echo-chamber collapse \\
PerfGD-analytic    & 0.36 & no settle (operator gradient) \\
PerfGD-learned     & 0.66 & no settle (negative result) \\
\textbf{PerfGD-structural} & \textbf{3.19} & \textbf{converges to realized optimum} \\
\bottomrule
\end{tabular}
\end{table}

\subsection{Lazy Deployment Stabilizes an Unstable Market (R6)}
\label{sec:lazy}
At a beyond-boundary configuration ($f=5$; exact-BR anchor $\hat m=1.205>1$
measured on the same seeds), the signed CRN $K$-step probe traces $\mu(K)$
for $K\in\{1,2,3,5,8,12,20\}$ (Figure~\ref{fig:lazy}). A one-parameter fit
of \eqref{eq:lazy} gives inner contraction $c=0.661$, and both
parameter-free predictions land. The deadbeat sign flip is predicted at
$\Kdb=1.46$, and measured medians run $+0.021$ at $K{=}1$ and $-0.064$ at
$K{=}2$. The stability-window exit is predicted at $\Kmax=5.75$, and
measured $|\mu|$ is $0.695$ at $K{=}5$ (inside) and $1.541$ at $K{=}8$
(outside). Laziness measurably keeps this RRM-unstable market stable for
$K\lesssim6$ and inherits the cobweb's divergence beyond, so retraining
\emph{cadence} is a stability control. The measured $\geff/\gamma$ falls
monotonically through $1$ between $K=5$ and $K=8$, the stiff branch of the
two-branch law. Beyond-boundary per-seed scatter is large by construction
and the medians carry the result. The clean tight-fit demonstration lives in
the contracting regime (measured $+0.78/+0.70/+0.37$ against fitted
$+0.88/+0.68/+0.36$ at $K=1/3/8$) and is locked in the test suite.

\begin{figure*}[t]
\centering
\includegraphics[width=\textwidth]{figures/lazy_deploy_gamma_eff.png}
\caption{Lazy deployment at a beyond-boundary configuration
($\hat m=1.205$). Left: the measured signed $K$-step map (medians over
seeds; faint dots are per-seed readings) against the one-parameter fit
$\mu(K)=-m+c^K(1+m)$ with $c=0.661$, its $K\to\infty$ asymptote $-m$, and
the predicted deadbeat cadence. Right: the implied effective curvature
$\geff(K)/\gamma$, predicted against measured, with the stiff branch
decaying through the exact-RRM line between $K=5$ and $K=8$.}
\label{fig:lazy}
\end{figure*}

\subsection{Scaling and Estimator Tuning}
\label{sec:tuning}

\begin{figure*}[t]
\centering
\includegraphics[width=\textwidth]{figures/universe_scaling.png}
\caption{Factor-model scaling (R5) with data-calibrated per-bond
dispersion. Left: $\rho(M)$ is flat in universe size from $d=8$ to $128$
bonds and pinned to the worst scalar modulus, far below the $\rho=1$
boundary (dotted). Right: the truncation bound $O(\lambda_{k+1}(C))$
against the measured error $|\rho-\rho_k|$ at $d=128$, showing three to four
orders of magnitude of slack at every retained-factor count $k$.}
\label{fig:universe}
\end{figure*}

\emph{Factor scaling (R5).} With per-bond volatilities dispersed at the
data-calibrated cross-sectional coefficient of variation, $\rho(M)$ is flat
at ${\approx}0.50$ from $d=8$ to $128$ bonds (Figure~\ref{fig:universe}) and
pinned to the worst scalar modulus. On this calibration, correlation does
not manufacture cross-sectional instability, and the fragile mode is
idiosyncratic. The Woodbury path computes $d=128$ in $0.12$\,s and matches
the dense eigensolve to machine precision; the truncation bound holds with
3--4 orders of magnitude of slack.

\emph{Sinkhorn blur (R1 instrument).} Against the exact 1-D quantile
$\Wone$, the debiased-divergence bias is U-shaped in the scale-relative blur
at a fixed iteration budget, with log-domain under-convergence below and
entropic over-blur above. Its minimum sits at $0.02\times$ std ($6.1\%$
ground-truth bias; $1.2\%$ on the config's CRN samples). We bake that in as
the default, and being scale-relative it transfers to real-unit
configurations unchanged.

\emph{Robust radius (R4).} Across normal, heavy-tailed, and skewed estimate
distributions the $z\,s$ radius over-covers ($0.99$--$1.00$ at $95\%$
nominal). Only contamination from railed-probe patterns degrades both radii
($0.938$), which is the honest $n_{\mathrm{req}}$ limit. On the actual CRN
probe estimates the quantile/$z\,s$ multiplier is $0.50$, so $z\,s$ is
binding and adequately calibrated. We keep the quantile-guarded radius for
suspected contamination.

%% ========================================================================
\section{Conclusion}
REFLEX turns the performative-prediction stability theorem from an
abstraction into an instrument. We compute $(\eps,\beta,\gamma)$ from
microstructure primitives, extend the boundary in closed form to correction,
competition, finite samples, dimension, and cadence, verify each extension
predict-then-verify against a learned loop, evaluate it as a daily fragility
index on 36 years of real market data, and machine-check the identities end
to end. Two findings matter beyond this model. Competition among retraining
dealers amplifies instability by the effective dealer count, which makes
certification of algorithmic liquidity provision a market-level rather than
a desk-level problem. And a learned corrected loop stabilizes where blind
retraining collapses only when its response model is anchored to market
structure. The free-form alternative fails with the same data and more
capacity, so a desk should spend its modelling budget on structural
anchoring rather than on network capacity.

These results hold under idealizing assumptions. The calibration is
proxy-level rather than trade-level TRACE: VIX-implied spreads, no
per-dealer inventories, a structurally scaled toxic channel, and a
degenerate crisis cell ($k=0$), so regime ordering is data-driven while
absolute critical gains are not. TRACE Enhanced calibration removes that
limit. At default-like constants the self-consistent fixed point never
destabilizes because dealers widen defensively, so measured crossings
describe the retraining map at the operating spread. Beyond the boundary,
probe readings are finite-difference diagnostics with seed-level
bifurcation, and medians with robust bands carry the results. Measured
moduli are protocol-dependent and comparable within, not across, protocols.
The structural loop's optimum is the \emph{realized} one, benchmarked
against independent structural fits rather than the frozen-reference closed
form, with the gap attributed to saturation, liquidity inflation, and
severity drift; its convergence is to the noise ball of an estimated
gradient. The free-form learned correction remains a negative result by
design. The Lean~4 skeletons are reviewed but not yet compiled, so the 66
numerical certificates are the verification of record. Future work:
trade-level calibration, inventory-state loop dynamics (state-dependent
performativity), richer policy classes, and compiling the formal layer.

\section*{Reproducibility}
\label{sec:repro}
%% DOUBLE-BLIND TOGGLE (step 2 of 2; step 1 is the \documentclass line).
%% To anonymize: comment out the \footnote below and uncomment the one
%% under it, after creating the anonymized mirror so the URL resolves.
The full framework is openly available:\footnote{%
\url{https://github.com/vignesh-nagarajan-vn/REFLEX/}}
%\footnote{Anonymized mirror for review: \url{https://anonymous.4open.science/r/REFLEX/}}
the six derivation documents D1--D6, the simulator, the operator, all four
loop modes, the estimator suite, the calibration pipeline with its data
catalogue and rejected-sources log, the 66-certificate verification layer,
the Lean sources, and the illustrated reports of both paper-grade runs.
All results in this paper are regenerated by one command
(\code{run\_all --profile full}, ${\sim}25$ CPU-minutes, deterministic from
configuration and seed).

%% ------------------------------------------------------------ references --
\bibliographystyle{ACM-Reference-Format}
\bibliography{references}

\end{document}
